\part{MyData}

\chapter{MyData --- Panoramica}
\label{ch:mydata}

MyData è il gestore risorse integrato di FairWindSK.
Fornisce un'interfaccia ottimizzata per il tocco per tutte le risorse nautiche archiviate
sul server Signal~K: waypoint, rotte, regioni, note, carte, tracce storiche e file.

Apri MyData toccando \iconlabel{database}{MyData} nella Barra Inferiore.

\screenshot{mydata-main-tabs}{MyData --- barra delle schede con le sette schede visibili}{}

\begin{notebox}
Tutte le risorse in MyData sono archiviate sul server Signal~K, non sul dispositivo
FairWindSK. Sono accessibili da qualsiasi dispositivo connesso allo stesso server
Signal~K simultaneamente.
\end{notebox}

\section{Panoramica delle Schede MyData}

MyData è organizzato in sette schede. Ogni scheda si carica in modo pigro --- viene
creata solo al primo tocco, mantenendo veloce l'avvio.

\rowcolors{2}{LtGray}{white}
\begin{longtable}{C{0.7cm}L{2.8cm}L{10.2cm}}
\toprule
\rowcolor{Navy}
\color{white}\sffamily\bfseries N. &
\color{white}\sffamily\bfseries Nome &
\color{white}\sffamily\bfseries Contenuto \\
\midrule
0 & Waypoint & Posizioni nominate: ancoraggio, ostacoli, marine, punti pesca, boe di regata. \\
1 & Rotte    & Sequenze ordinate di waypoint che formano traversate pianificate. \\
2 & Regioni  & Aree geografiche di confine --- zone di ricerca SAR, zone di esclusione, limiti di ancoraggio. \\
3 & Note     & Annotazioni di testo libero associate a posizioni. \\
4 & Carte    & Catalogo sorgenti tile nautiche per i fornitori di carte installati. \\
5 & Tracce   & Traccia GPS storica --- la sequenza registrata delle posizioni dell'imbarcazione. \\
6 & File     & Browser diretto del filesystem per i file accessibili tramite Signal~K. \\
\bottomrule
\end{longtable}

\section{La Barra degli Strumenti Comune}

Le schede 0--4 (Waypoint, Rotte, Regioni, Note, Carte) condividono un layout comune:
una tabella scorrevole con le risorse dal server Signal~K e una barra degli strumenti
con i pulsanti icona:

\rowcolors{2}{LtGray}{white}
\begin{longtable}{L{2.8cm}L{10.9cm}}
\toprule
\rowcolor{Navy}
\color{white}\sffamily\bfseries Pulsante &
\color{white}\sffamily\bfseries Azione \\
\midrule
\iconlg{arrow-right-google}~Apri & Apre la risorsa selezionata --- mostra i dettagli o l'anteprima. \\
\iconlg{edit-google}~Modifica     & Apre il cassetto editor per la risorsa selezionata. \\
\iconlg{widget-add-google}~Nuovo  & Apre il cassetto editor per creare una nuova risorsa. \\
\iconlg{delete-google}~Elimina    & Elimina la risorsa selezionata dopo conferma. Irreversibile. \\
\iconlg{route-import-iec}~Importa & Importa da file GeoJSON o JSON. \\
\iconlg{file-export-google}~Esporta & Esporta in file GeoJSON. \\
\iconlg{refresh-google}~Aggiorna  & Forza il ricaricamento di tutte le risorse dal server Signal~K. \\
\bottomrule
\end{longtable}

\chapter{Waypoint}
\label{ch:waypoint}

I waypoint sono posizioni nominate archiviate sul server Signal~K: ancoraggi, ostacoli,
marine, punti di rifornimento, punti pesca, boe di regata e qualsiasi altra posizione
che si voglia registrare e verso cui navigare.

\screenshot{mydata-waypoints-list}{MyData --- scheda Waypoint, vista elenco}{}

\section{Colonne della Tabella Waypoint}

Nome, Descrizione, Tipo (es.\ \code{alarm-mob} per i waypoint POB), Latitudine, Longitudine,
Timestamp.

\section{Creazione di un Waypoint}

Tocca \key{Nuovo}. Si apre il cassetto Editor Waypoint con i campi: Nome, Descrizione, Tipo,
Latitudine, Longitudine, Altitudine.

\section{Navigare verso un Waypoint}

Dalla vista dettaglio, tocca \key{Naviga}.
FairWindSK imposta il waypoint come destinazione di navigazione attiva su Signal~K.
Immediatamente i widget BTW, DTG, TTG ed ETA nella Barra Superiore mostrano valori
relativi a questo waypoint.

\chapter{Rotte}
\label{ch:rotte}

Le rotte sono sequenze ordinate di waypoint che formano traversate pianificate.

\screenshot{mydata-routes-list}{MyData --- scheda Rotte, vista elenco}{}

\section{Colonne della Tabella Rotte}

Nome, Descrizione, Geometria (es.\ \code{LineString}, \code{sar-spiral}, \code{sar-parallel}),
Punti, Timestamp.

\begin{tipbox}
I percorsi di ricerca SAR creati dalla funzione POB appaiono automaticamente nella
scheda Rotte. Hanno tipi di geometria \code{sar-spiral} e \code{sar-parallel} e possono
essere esportati come GeoJSON per i registri di viaggio.
\end{tipbox}

\chapter{Regioni}
\label{ch:regioni}

Le regioni sono aree di confine geografiche definite come poligoni GeoJSON.
Vengono usate per aree di ricerca SAR (create automaticamente dalla funzione POB),
zone marine protette, zone di esclusione, limiti di ancoraggio e zone di pesca.

\section{Colonne della Tabella Regioni}

Nome, Descrizione, Geometria, Timestamp.

\chapter{Note}
\label{ch:note}

Le note sono annotazioni di testo libero associate a posizioni geografiche.
Funzionano come post-it sulla carta --- utili per registrare consigli di pilotaggio,
ostacoli locali non in carta, dettagli di marine o osservazioni dell'equipaggio.

\screenshot{mydata-notes-list}{MyData --- scheda Note, vista elenco con titolo, descrizione, href, tipo MIME}{}

\section{Colonne della Tabella Note}

Titolo, Descrizione (per i file di testo importati: i primi 120 caratteri del contenuto),
Href (percorso file o URL dell'allegato), Tipo MIME, Timestamp.

\section{Importazione Note}

La scheda Note accetta un insieme più ampio di tipi di file di importazione:
file GeoJSON/JSON, file di testo (\code{.txt}, \code{.md}, \code{.html}), immagini
(\code{.jpg}, \code{.jpeg}, \code{.png}) e file PDF.
Per i file non-GeoJSON, FairWindSK crea automaticamente una nota con il nome del file
come titolo e l'href impostato sul percorso assoluto del file.

\chapter{Carte}
\label{ch:carte}

La scheda Carte elenca le sorgenti di tile nautici disponibili registrate sul server Signal~K.

\screenshot{mydata-charts-list}{MyData --- scheda Carte, vista elenco}{}

\section{Colonne della Tabella Carte}

Nome, Descrizione, Formato (\code{mbtiles}, \code{kap}, \code{gemf},
\code{tilemapresource}, \code{package}), Identificatore, URL Carta, Timestamp.

\chapter{Tracce --- Storico}
\label{ch:tracce}

La scheda Tracce mostra la traccia GPS storica dell'imbarcazione.

\screenshot{mydata-tracks-list}{MyData --- scheda Tracce, tabella dei punti traccia con ora, lat/lon, altitudine}{}

\section{Colonne della Tabella Tracce}

Ora (timestamp formattato in ora locale), Latitudine, Longitudine, Altitudine.

\section{Esportazione della Traccia}

Tocca \key{Esporta} per salvare la traccia completa come GeoJSON \code{LineString}.
Il nome file predefinito è \code{tracks-history.geojson}.

\begin{tipbox}
Esporta la traccia dopo ogni traversata e archivial con il giornale di bordo.
Il file GeoJSON può essere importato nella maggior parte degli strumenti GIS,
plotter e software di pianificazione di viaggio.
\end{tipbox}

\chapter{File}
\label{ch:file}

La scheda File è un browser diretto del filesystem per i file e le directory
accessibili sul dispositivo che esegue FairWindSK.

\screenshot{mydata-files-browser}{MyData --- scheda File, browser directory con barra degli strumenti}{}

\section{Navigazione e Ricerca}

Naviga tra le cartelle toccando due volte le righe directory.
Usa il pulsante \iconlg{home}~Home per tornare alla home directory.
Usa \iconlg{arrow-up-google}~Su per salire di un livello.

Per la ricerca, tocca \iconlg{search}~Cerca. Le opzioni includono:
\textbf{Distinzione maiuscole/minuscole}, \textbf{Cerca nascosti} e \textbf{Cerca di sistema}.

\section{Azioni sui File}

\rowcolors{2}{LtGray}{white}
\begin{longtable}{L{3.6cm}L{10.1cm}}
\toprule
\rowcolor{Navy}
\color{white}\sffamily\bfseries Azione &
\color{white}\sffamily\bfseries Descrizione \\
\midrule
\iconlg{arrow-right-google}~Apri &
  Immagini: visore immagini. PDF: visore PDF. File di testo: visore testo. \\
\iconlg{widget-add-google}~Nuova cartella &
  Crea una nuova cartella nella directory corrente. \\
\iconlg{content-copy-google}~Copia &
  Segna gli elementi selezionati per la copia. \\
\iconlg{content-cut-google}~Taglia &
  Segna gli elementi selezionati per lo spostamento. \\
\iconlg{content-paste-google}~Incolla &
  Incolla nella cartella corrente. \\
\iconlg{delete-google}~Elimina &
  Elimina gli elementi selezionati dopo conferma. Permanente --- nessun cestino. \\
\bottomrule
\end{longtable}

\begin{warningbox}
L'eliminazione dei file è permanente.
Non esiste cestino né funzione annulla per le operazioni sui file.
Prestare particolare attenzione all'eliminazione di directory --- tutto il contenuto
viene eliminato ricorsivamente.
\end{warningbox}
